\section{Geometria sferyczna} % lang-pl
\label{sekcja_sferyczna}

Druga geometria nieeuklidesowa jest starsza od pierwszej, tylko nikt jeszcze nie zauważy, że to geometria nieeuklidesowa: będzie to zwykła geometria na kuli. % lang-pl
W tej sekcji zobaczymy, że spełnia ona hipotezę kąta rozwartego, a potem poznamy jedno z jej twierdzeń, które przez sto lat zajmie największych geometrów. % lang-pl

Na sferze „prostymi'' są koła wielkie, dwie z nich zawsze się przecinają, a suma kątów trójkąta jest większa od $\pi$; Lambert zauważy tę analogię z hipotezą kąta rozwartego \cite[s. 286]{eves1_1972} (zobacz sekcję \ref{sekcja_lambert}). % lang-pl
Pole trójkąta sferycznego jest proporcjonalne do jego nadmiaru: dla sfery o promieniu $r$ i nadmiarze $E$ w stopniach jest ono równe $\pi r^2 E/180^\circ$ (Eves \cite[s. 289--290, zad. 10]{eves1_1972}). % lang-pl

\begin{theorem}
[Girarda] % lang-pl
\index{twierdzenie!Girarda} % lang-pl
	Pole trójkąta sferycznego o kątach $\alpha$, $\beta$, $\gamma$ na sferze o promieniu $R$ jest równe $R^2(\alpha + \beta + \gamma - \pi)$. % lang-pl
\end{theorem}

Albert Girard opublikuje je w 1629 roku; wcześniejszy, niepublikowany dowód znajdzie się w rękopisach Thomasa Harriota z 1603 roku. % lang-pl
Odpowiednikiem w geometrii hiperbolicznej będzie pole równe defektowi (zobacz sekcję \ref{sekcja_suma_katow}). % lang-pl
\index[persons]{Girard, Albert}%
\index[persons]{Harriot, Thomas}%

Eves \cite[s. 314--315]{eves1_1972} mówi krótko, że hipotezie kąta rozwartego odpowiada druga geometria nieeuklidesowa, że jest ona w pewnym sensie bardziej skomplikowana od geometrii Łobaczewskiego i że nie rozwija jej w swojej książce. % lang-pl
\todofoot{Geometria eliptyczna (sfera z utożsamionymi punktami antypodalnymi; nazwę wprowadzi Klein w 1871) oraz trygonometria sferyczna -- potrzebne źródło poza Evesem.} % lang-pl

% Sferyczna osobno?
\begin{proposition}
[twierdzenie Lexella] % lang-pl
	Wierzchołki trójkątów sferycznych o ustalonej podstawie $AB$ i polu powierzchni leżą na okręgu małym (powstałym przez przecięcie sfery płaszczyzną) przechodzącym przez punkty antypodalne do wierzchołków $A$, $B$ z podstawy. % lang-pl
\index{twierdzenie!Lexella} % lang-pl
\end{proposition}

Anders Johan Lexell napisze około 1777 roku pracę na temat tego twierdzenia, razem z dwoma dowodami: trygonometrycznym i geometrycznym (wydrukowana będzie w 1784 roku). % lang-pl
Inni podadzą inne dowody: % lang-pl
Leonhard Euler w~1778 (kolega Lexella; jego dowód zostanie wydrukowany dopiero w 1797 roku), % lang-pl
Adrien-Marie Legendre w~1800, % lang-pl
Jakob Steiner w~1827, % lang-pl
Carl Friedrich Gauss w~1841, % lang-pl
Paul Serret w~1855, % lang-pl
Joseph-Émile Barbier w~1864. % lang-pl
\index[persons]{Lexell, Anders Johan}%
\index[persons]{Serret, Paul}%
\index[persons]{Barbier, Joseph-Émile}%
\index[persons]{Steiner, Jakob}%

Twierdzenie Lexella jest sferycznym odpowiednikiem twierdzeń I.37 i I.39 z \emph{Elementów}: trójkąty płaskie o wspólnej podstawie i równych polach mają wierzchołki na prostej równoległej do podstawy. % lang-pl

% Uwaga: dane o Lexellu pochodzą z Wikipedii (Lexell's theorem); Eves twierdzenia nie omawia.
% https://en.wikipedia.org/wiki/Spherical_trigonometry => Girard
% Źródła: Eves s. 287, 306--307, 352 (zad. 6--7); Wikipedia: Eugenio Beltrami, Arthur Cayley, Felix Klein, Beltrami–Klein model, Non-Euclidean geometry.
% https://en.wikipedia.org/wiki/Beltrami%E2%80%93Klein_model
% na podstawie https://en.wikipedia.org/wiki/Beltrami–Klein_model
% model Kleina - Delta 2012-06

